% ============================================================
% tab/matriz_revision.tex
%
% COLUMNAS OBLIGATORIAS (enunciado): referencia, problema,
% método, datos o entorno, hallazgo principal, limitación.
% No sustituir, no reordenar. No copiar abstracts.
%
% CONTENIDO: solo los 5 trabajos principales (P1-P5). Los
% opcionales y las fuentes de actualidad se discuten en el texto.
%
% REGLA DE tabularx: 5 columnas W, pesos suman 5.00.
% ============================================================

\begin{table*}[!tp]
\centering
\caption{Matriz de revisi\'on de los trabajos principales. Las cifras est\'an
transcritas de la fuente. Las tareas, conjuntos y sistemas evaluados
difieren, por lo que la comparaci\'on directa requiere cautela.}
\label{tab:matriz}
\footnotesize
\setlength{\tabcolsep}{3.5pt}
\begin{tabularx}{\textwidth}{@{}
    P{1.9cm}
    >{\centering\arraybackslash}p{0.6cm}
    W{0.95}
    W{0.95}
    W{0.85}
    W{1.15}
    W{1.10}
  @{}}
\toprule
\textbf{Referencia} & \textbf{A\~no} & \textbf{Problema} &
\textbf{M\'etodo} & \textbf{Datos / entorno} &
\textbf{Hallazgo principal} & \textbf{Limitaci\'on}\\
\midrule

Grother et al.\ \cite{grother_nistir8280_2019} & 2019 &
Diferencias de exactitud entre grupos demogr\'aficos en verificaci\'on (1:1)
e identificaci\'on (1:N) &
Evaluaci\'on independiente a umbral fijo, reportando por separado falsos
positivos y falsos negativos &
18.27\,M de im\'agenes de 8.49\,M de personas (fichas policiales, visas,
inmigraci\'on y cruces fronterizos de EE.\,UU.); 189 algoritmos de 99
desarrolladores &
Los falsos positivos var\'ian entre 10 y m\'as de 100 veces entre grupos y
son mayores en mujeres, ni\~nos y adultos mayores; los falsos negativos
var\'ian usualmente menos de 3 veces y dependen de la calidad de imagen &
No explica causas; no usa im\'agenes de internet ni de videovigilancia;
algoritmos prototipo evaluados tal como fueron enviados
\\
\addlinespace[3pt]

Kotwal y Marcel \cite{kotwal_fairness_2026} & 2026 &
Disparidades persistentes en modelos de reconocimiento facial modernos &
Revisi\'on que organiza causas, conjuntos de datos, m\'etricas y
mitigaciones (pre, intra y posprocesamiento) &
Literatura publicada hasta 2025; evaluaciones del NIST y norma ISO/IEC
19795-10 &
Las disparidades aparecen en todo el rango de umbrales; combinan datos de
entrenamiento, tono de piel, calidad de imagen y factores no demogr\'aficos;
balancear los datos no las elimina &
Centrada en raza y sexo, con la edad tratada de forma marginal; sin
experimentos propios; poca evidencia en baja resoluci\'on, t\'ipica de la
vigilancia
\\
\addlinespace[3pt]

Vakhshiteh et al.\ \cite{vakhshiteh_adversarial_2021} & 2021 &
Vulnerabilidad del reconocimiento facial profundo a ejemplos adversarios &
Revisi\'on con taxonom\'ia de ataques (cuatro orientaciones) y de defensas
(tres estrategias) &
Estudios sobre FaceNet, ArcFace, SphereFace y servicios comerciales, con
conjuntos como LFW y CASIA-WebFace &
Anteojos impresos, gorras con LED infrarrojos, luz proyectada o parches
enga\~nan a sistemas reales; las defensas cubren ataques espec\'ificos sin
generalizar &
Limitada a ataques adversarios; los estudios revisados usan condiciones de
captura controladas; sin comparaci\'on cuantitativa com\'un
\\
\addlinespace[3pt]

Fussey et al.\ \cite{fussey_assisted_2020} & 2021 &
Operaci\'on real del reconocimiento facial en vivo por la polic\'ia y papel
de la discreci\'on humana &
Etnograf\'ia (observaci\'on y entrevistas) y an\'alisis de alertas y
bit\'acoras &
Pilotos de South Wales Police (35\,h observadas) y Polic\'ia Metropolitana
de Londres (m\'as de 50\,h), 2016--2019 &
En Londres, de 43 alertas se juzgaron cre\'ibles 27 y solo 9 fueron
correctas; el umbral, las listas de vigilancia y la discreci\'on del
operador determinan el resultado &
Solo dos fuerzas policiales del Reino Unido; muestras peque\~nas; muchas
alertas nunca verificadas; no mide sesgo demogr\'afico
\\
\addlinespace[3pt]

Almeida et al.\ \cite{almeida_ethics_2022} & 2022 &
Rendici\'on de cuentas del uso policial del reconocimiento facial &
An\'alisis comparativo de marcos regulatorios y jurisprudencia &
EE.\,UU., Uni\'on Europea y Reino Unido; casos judiciales y sanciones de
autoridades de datos &
El reglamento europeo exige evaluaciones de impacto y trata la biometr\'ia
como categor\'ia especial; EE.\,UU. carece de ley federal; propone diez
preguntas para regular su uso &
Discursivo, sin evaluaci\'on emp\'irica; no cubre Am\'erica Latina;
anterior a la aprobaci\'on del Reglamento Europeo de IA
\\

\bottomrule
\end{tabularx}
\end{table*}